% =====================================================================
\section{Stage 1.1 (v5.1): Three-Feature Expansion}
\label{sec:stage11-v51}

After v5.0.1 fixed the Faith bias, three orthogonal features were
added to v5.1, each addressing a separate dimension of the
specification originally laid out at the start of this Part:

\begin{enumerate}
\item \textbf{Ahistorical event catalog} (Requirement~7).
  Twenty events spanning physical anomalies, ideological breaks,
  resource shocks, climate collapses, and global disruptions, with
  per-event year-windows and per-year hazard probabilities.
\item \textbf{Frequency modes} (Requirement~8).
  Three operational settings —~\emph{sporadic} (multiplier
  $\times 1$), \emph{frequent} ($\times 3$), and \emph{sparse}
  ($\times 0.3$)~— acting as a global multiplier on event hazard
  rates. This permits ``crisis-rich'' and ``crisis-quiet'' worlds to
  be compared with the same mechanics.
\item \textbf{Encounter mechanism} (Requirement~4).
  Four channels (\emph{biological},
  \emph{technological}, \emph{ideological}, \emph{cosmological}) ×
  four intensities (\emph{soft}, \emph{moderate}, \emph{intense},
  \emph{catastrophic}), with intensity evolving generation by
  generation after first contact. This realises the ``unknown
  encounter'' world type.
\end{enumerate}

The result was a single simulator (\texttt{world\_sim\_v5\_1.py},
1212 lines) that handles all five world families plus the
encounter channel within one consistent action--state loop.

% =====================================================================
\section{Stage 1.2 (v5.2--v5.4): Spotlight, Time Extension, Multi-Civ}
\label{sec:stage12-v52-v54}

Three further v5.x increments addressed the ``protagonist'' problem
(Requirement~10) and laid groundwork for multi-civilisation
parallelism (Requirement~11):

\paragraph{v5.2 (Spotlight, three-tier).}
For each simulation run, one \emph{individual} is followed across
generations with a per-decision chronicle (eight emotion types,
crisis responses, career transitions, ahistorical encounters). One
\emph{family} (four branches: main line, two cadet branches, one
adoptee/affine line) is tracked with an ASCII family tree and a
shared lineage chronicle. The full \emph{civilisation} is summarised
per generation as a six-dimensional state snapshot
(family-knowledge, education, institution, trade, assets, urban).
This is the v5.2 ``three-tier Spotlight'': individual / family /
civilisation observed simultaneously within one simulation.

\paragraph{v5.3 (3000-year horizon).}
The era schedule was extended past 2021 with two future eras —
\emph{Future\_AI\_Era} (2021--2500) and \emph{Future\_Posthuman}
(2500--3000) — and a complementary catalog of ten future
ahistorical events (\texttt{AI\_Singularity},
\texttt{Climate\_Tipping}, \texttt{Posthuman\_Transition},
\texttt{Resource\_Depletion}, etc.). The simulation now spans 75
generations (3000 years) by default rather than the historical 50
generations (2000 years).

\paragraph{v5.4 (Multi-civ parallel).}
Four civilisations (Japan, China, Europe, Islamic) are run in
parallel with a small inter-civ interaction catalog (six pairs ×
four channels: trade, war, knowledge\_diffusion,
disease\_exchange). The era-dependent interaction catalog encodes,
for example, Japan--China at year 0--700 with high knowledge
transmission probability (Buddhism, kanji), Europe--Islamic at
1100--1300 with both war (Crusades) and knowledge transmission
(translation movement). This was the first concrete realisation of
Requirement~11 ``multi-civilisation parallel'' even at small
scale.

% =====================================================================
\section{Stage 2 (v5.5): Cohort Spotlight}
\label{sec:stage2-v55}

The v5.2 Spotlight follows one named individual; v5.5 expands this
to a \textbf{cohort of 100+ named individuals} chronicled
simultaneously, addressing the original Stage 2 milestone of the
Telescope Architecture roadmap (Section~\ref{sec:vision-roadmap}).

\paragraph{Diversity grid sampling.}
The cohort is selected by stratified sampling across a
three-dimensional grid: (decision strategy) ×
(initial region) × (founding occupation tier $\in$ \{agrarian,
notable, craft-trade, urban\}). For 11 strategies × 6 regions × 4
tiers = 264 cells, with 100 cohort slots, roughly 100/264 cells get
filled. This guarantees the cohort is not 100 NRMO farmers in
North Kyushu, but spans strategies and starting positions.

\paragraph{Five highlight categories.}
For each run, the cohort is automatically classified into:
\begin{itemize}
\item \emph{Most successful} (highest final education + assets);
\item \emph{Most resilient} (most crises survived);
\item \emph{Most dynamic} (most career transitions);
\item \emph{Tragic loss} (lineage absorbed);
\item \emph{Stayed humble} (remained agrarian throughout).
\end{itemize}
Each category produces a top-five highlight list with brief
narrative summary, allowing the human reader to rapidly identify
illustrative cases without reading all 100 chronicles.

\paragraph{Cohort emotional aggregation.}
Per-generation emotion records across all cohort members are
aggregated into a cohort-wide emotional arc histogram. In a typical
Normal-Earth, sporadic-frequency, 3000-year run this resolves to
roughly: hope $\sim$42\%, equanimity $\sim$26\%, loss $\sim$11\%,
disorientation $\sim$8\%, fear $\sim$7\%, anxiety $\sim$5\%, grief
$\sim$1\%. This is the population-level emotional shape of a 3000
year simulation, observed at individual resolution.

% =====================================================================
\section{Stage 3 (v6.0): Cultural Modules and Inter-Civ Interaction}
\label{sec:stage3-v60}

Stage 3 of the roadmap (10$^7$ agents, 10+ Cultural Modules,
inter-civ interaction) was attempted in v6.0 with the following
results:

\paragraph{Nine Cultural Modules.}
Beyond the four v5.4 modules, five additional civilisations were
modelled:
\begin{itemize}
\item \emph{Indic}: Mauryan -- Gupta -- Sultanate -- Mughal --
  British -- Independent India eras; joint family
  \texttt{distribution\_boost}~$=1.25$; mixed Hindu / Buddhist /
  Jain / Islamic Faith distribution.
\item \emph{SubSaharan}: Iron Age kingdoms -- Trans-Saharan --
  Atlantic Slave Trade -- Colonial Scramble -- Postcolonial;
  communal extended family inheritance.
\item \emph{Polynesian}: Lapita expansion -- Classical --
  Contact Era (1769+, Cook contact) -- Colonial -- Modern Pacific;
  primogeniture-with-genealogy emphasis.
\item \emph{Steppe Nomadic}: Scythian -- Turkic Khanates --
  Mongol Empire (1200--1400 peak) -- Post-Mongol -- Russian/Qing
  -- Modern Sedentary; high militant Faith share (warrior tradition).
\item \emph{IndigenousAmericas}: Classic Civilisation
  (Maya, Teotihuacan) -- Postclassic (Aztec, Inca) -- Conquest
  Catastrophe (1521--1700, $\sim$90\% population loss per Cook
  1998) -- Colonial Mestizo -- Independence -- Modern Indigenous.
\end{itemize}

\paragraph{Eighteen interaction pairs.}
The interaction catalog grew from six to eighteen civilisation
pairs, encoding, e.g., China-Indic Buddhist transmission and
Xuanzang's journey, China-Steppe Mongol conquest of 1200--1400
(intensity 0.95), Europe-IndigenousAmericas Conquest of 1492
(disease intensity 1.0, war intensity 0.95), Europe-Polynesian
Cook-era contact (disease intensity 0.95), Islamic-SubSaharan
Trans-Saharan trade and Islamic transmission, Indic-Islamic
Sultanate and Mughal eras.

\paragraph{Scaling: a partial achievement.}
Stage 3's nominal target was 10$^7$ agents. In a 4\,GB-RAM,
non-numba environment we reached 9 civilisations × 100{,}000 agents
= 900{,}000 agents (= $\approx 10^6$) in 66 seconds, with 886
inter-civ interactions over 75 generations. The 10$^7$ target was
\emph{not} achieved; numba JIT installation failed in the
constrained environment, and pure-numpy memory pressure (estimated
$\sim$6\,GB at 9 × 10$^6$) exceeded available RAM. The honest
status is therefore: Stage 3 \emph{partial} (1/10 of nominal scale,
all other features in place).

% =====================================================================
\section{v6.1: Unified Integrated Simulator}
\label{sec:v61-unified}

The v5.5 cohort Spotlight and the v6.0 multi-civ machinery were
integrated into a single configurable simulator,
\texttt{world\_sim\_v6\_1\_unified.py}, with four scale presets:

\begin{center}
\begin{tabular}{lrrrrr}
\toprule
Preset & Per-civ agents & Total agents & Cohort/civ & Runtime & RAM \\
\midrule
\texttt{small}  &   5{,}000 &   45{,}000 &  50 &   $\sim$5\,s   & $<$0.5\,GB \\
\texttt{medium} &  50{,}000 &  450{,}000 & 100 &   $\sim$30\,s  & $\sim$1\,GB \\
\texttt{large}  & 200{,}000 & 1{,}800{,}000 & 200 &  $\sim$170\,s & $\sim$1.2\,GB \\
\texttt{custom} & user-defined & --- & --- & --- & --- \\
\bottomrule
\end{tabular}
\end{center}

A single run produces, per civilisation: a cohort summary
markdown, a family chronicle markdown, and a per-generation
trajectory CSV; plus cross-civilisation comparisons
(\texttt{cross\_civ\_strategy.csv}), an inter-civ interaction
log (\texttt{interaction\_summary.md}), an executive overall report,
and a full data JSON dump. \textbf{Five tiers are observed
simultaneously in one simulation:} individual (cohort) / family
(four branches per civ) / civilisation (nine modules) / inter-civ
(eighteen pairs) / world (random + ahistorical events).

% =====================================================================
\section{Stage 4 Partial (v6.2): Memetic, Black Swan, Counterfactual}
\label{sec:stage4-v62}

Three of the four Stage 4 features were implemented in v6.2; the
fourth (GPU acceleration to 10$^8$ agents) remained
environment-blocked. The three implemented features are exposed as
optional CLI flags so the v6.1 baseline behaviour is preserved by
default.

\subsection{Memetic Dynamics}
\label{subsec:memetic}

Strategy distributions are no longer fixed. At each generation
end, per-strategy fitness is computed as
\begin{equation}
\textsc{Fitness}_s
= \frac{\#\text{alive}_s}{\#\text{total}_s}
\times \left(\tfrac{1}{2} +
\tfrac{1}{2}\,(\overline{\text{edu}}_s + \overline{\text{assets}}_s)\right),
\end{equation}
where the first factor is per-strategy survival and the second is
mean prosperity among alive members. A replicator step then nudges
the strategy distribution toward higher-fitness strategies, with a
drift rate (default 0.05) controlling adoption speed. Newly-reset
collateral lineages are reassigned strategies according to the
updated distribution.

This realises memetic evolution: ``Faith\_Militant'' that
self-undermines (the Stage 1 empirical results above) sees declining
share over many generations; ``NRMO\_vNext'' steady-state share
rises in volatile worlds.

\subsection{Black Swan Events}
\label{subsec:black-swan}

A separate catalog of nine ultra-rare ($p \in [10^{-5}, 10^{-4}]$
per year) events with per-event scope (global / regional /
civ-pair) and per-event mortality probability:

\begin{center}
\small
\begin{tabular}{lllrr}
\toprule
Event & Period & Scope & Shock & Mortality \\
\midrule
Volcanic\_Winter\_Year       & 0--3000  & global    & 0.50 & 10\% \\
Megaplague\_Pandemic         & 0--3000  & regional  & 0.40 & 20\% \\
Great\_Solar\_Flare\_Carrington & 1850--3000 & global & 0.30 & 5\% \\
Asteroid\_Regional\_Impact   & 0--3000  & regional  & 0.70 & 30\% \\
Multi\_Empire\_Collapse      & 0--2021  & civ\_pair & 0.55 & 15\% \\
Climate\_Tipping\_Cascade    & 2050--3000 & global  & 0.45 & 10\% \\
AI\_Misalignment\_Crisis     & 2050--3000 & global  & 0.60 & 5\% \\
Antibiotic\_Apocalypse       & 2030--3000 & global  & 0.35 & 18\% \\
Posthuman\_Schism            & 2500--3000 & global  & 0.50 & 5\% \\
\bottomrule
\end{tabular}
\end{center}

An intensity multiplier permits stress testing: with multiplier
$\times 50$, several events fire per typical 3000-year run.

\subsection{Counterfactual Mode}
\label{subsec:counterfactual}

Six pre-built counterfactual scenarios:

\begin{center}
\small
\begin{tabular}{ll}
\toprule
Scenario & Description \\
\midrule
\texttt{no\_black\_death}  & Suppress Black Death pandemic 1347--1353 \\
\texttt{no\_conquest}      & IndigenousAmericas not conquered (1492 contact peaceful) \\
\texttt{no\_cook}          & Polynesia stays isolated (no European contact) \\
\texttt{no\_mongol}        & No Mongol Empire (1200--1400 conquests prevented) \\
\texttt{no\_industrial}    & Industrial Revolution suppressed \\
\texttt{early\_islam}      & Islamic civilisation rises 200 years earlier \\
\bottomrule
\end{tabular}
\end{center}

Override types: \texttt{suppress\_event} (filter ahistorical
labels), \texttt{civ\_protection} (reduce incoming shock to a
civ in a year window), \texttt{civ\_extinction} (force absorption),
\texttt{suppress\_interaction} (block specific civ pair
interactions in a year window).

\paragraph{Empirical effect (single seed, small scale).}
At seed=100, scale=small, comparing baseline vs counterfactual:
IndigenousAmericas continuation 14.8\% $\to$ 16.8\% with
\texttt{no\_conquest} (+2.0\,pp); Polynesian 55.8\% $\to$ 58.4\%
with \texttt{no\_cook} (+2.6\,pp). The effects are modest because
the present implementation only modulates inter-civ shocks; the
era-fixed \texttt{base\_failure} inside the Cultural Module is
\emph{not} overridden. A fully era-dynamic counterfactual would
require deeper hooks and is deferred.

% =====================================================================
\section{Calibration (v6.2)}
\label{sec:calibration-v62}

Cultural Module parameters were re-tuned against historical
demography literature (Cook~1998; McEvedy \& Jones~1978;
Diamond~1997; Reich~2018; Cavalli-Sforza~2000). Each civilisation
was assigned a target range for 2000-year (Y0--Y2021) lineage
continuation probability, summarised below.

\begin{center}
\begin{tabular}{lrrrr}
\toprule
Civ & Pre-cal. & Post-cal. & Target & Status \\
\midrule
Japan & 0.87 & 0.87 & 0.78--0.92 & in range \\
China & 0.84 & 0.84 & 0.78--0.90 & in range \\
Europe & \textbf{0.29} & 0.46 & 0.40--0.70 & moved into range \\
Islamic & 0.84 & 0.84 & 0.70--0.88 & in range \\
Indic & 0.88 & 0.87 & 0.75--0.90 & in range \\
SubSaharan & 0.86 & 0.82 & 0.55--0.82 & moved into range \\
Polynesian & \textbf{0.84} & 0.51 & 0.40--0.72 & moved into range \\
Steppe & 0.77 & 0.77 & 0.50--0.78 & in range \\
IndigenousAmericas & \textbf{0.80} & 0.15 & 0.05--0.25 & moved into range \\
\bottomrule
\end{tabular}
\end{center}

The most consequential calibration was \texttt{IndigenousAmericas}
\texttt{Conquest\_Catastrophe} era \texttt{base\_failure}, raised
from $0.45$ to $0.65$ to reflect the $\sim$90\% population loss
1521--1620 documented by Cook~1998. With this single change the
simulated continuation rate fell from 80\% to 15\%, well within
the 5--25\% target band.

Parallel adjustments: Polynesian Contact-Era \texttt{base\_failure}
$0.20 \to 0.35$ (Marquesas 90\% loss);
Atlantic Slave Trade \texttt{base\_failure} $0.20 \to 0.30$
(forced demographic outflow); Europe Migration Period
\texttt{base\_failure} $0.18 \to 0.16$ (slightly relaxed).

The result is the v6.2 default profile: nine civilisations all
within target ranges under the seed=20260507 baseline. Calibration
remains seed-sensitive — Black Swan-enabled runs may push Europe
or Polynesian below the target range — but the central tendency is
empirically anchored.

% =====================================================================
\section{Honest Limits and Stage 4 / Stage 5 Outlook}
\label{sec:v62-limits-outlook}

\paragraph{What v6.2 has achieved.}
Stage~1 complete; Stage~2 complete (cohort of 100+ subjects
chronicled per civilisation); Stage~3 \emph{partial} (9 Cultural
Modules and 18 interaction pairs implemented at 10$^6$-agent scale,
not the 10$^7$ nominal target); Stage~4 \emph{partial} (Memetic,
Black Swan, Counterfactual implemented; GPU acceleration not).
Calibration to historical demography targets has been performed.

\paragraph{What v6.2 has \emph{not} achieved.}
The 10$^7$-agent target of Stage~3 awaits numba/CUDA capability.
Stage~4's GPU acceleration to 10$^8$ agents likewise.
Counterfactual override of era-fixed civilisation parameters is
implemented only via inter-civ shock modulation, not full
era-dynamic re-parameterisation. Memetic dynamics use
fitness-proportional drift but do not yet model strategy
\emph{transmission} between civs (cross-cultural meme spread).

\paragraph{Stage 5.}
Stage~5's full Telescope across 10+ Cultural Modules and 10$^9$
agents remains pending and is, in present hardware terms,
fundamentally GPU-constrained. The architecture from
Section~\ref{sec:vision-telescope} accommodates it; only
implementation hardware differs.

\paragraph{Source-fidelity claim.}
Throughout v5.0--v6.2 development, the architectural invariant of
governance-execution separation
(the governance--execution separation invariant) was preserved without
exception. NRMO defines the admissibility set $A_t$; the engine
$\Omega$ searches only inside it; the warn flag from
TimeHorizonLayer is never passed into NRMO\_CORE\@. No simulator
revision relaxed this discipline, and no shortcut to performance
came at its expense.

% =====================================================================
\section{Connections to Other Parts}

\begin{itemize}
\item This Part is the long-form specification of the simulation
ambition behind Part XI's
historical validation (the Part XI historical validation chapter).

\item The decision theories used here are the same as Part XI's strategy
set, plus \texttt{Faith} (vision-introduced).

\item The Telescope Architecture extends the multi-horizon principle of
Part II's TimeHorizonLayer: the
same single decision is evaluated at multiple resolution scales,
analogous to TimeHorizonLayer evaluating a decision at multiple time
horizons.

\item The Cultural Modules realise the open item from Part XII
(the Part XII methodology limits section: ``Multi-agent
inter-civilisation interaction''), now empirically calibrated against
historical demography literature
(Section~\ref{sec:calibration-v62}).

\item The Counterfactual Mode realises the long-deferred goal of doing
counterfactual analysis on NRMO's value, complementing the historical
8-scenario validation in Part XI; it is now operational with six
named scenarios (Section~\ref{subsec:counterfactual}).

\item The Individual Layer re-introduces the cognitive 8-tuple
Situation Parameter Model from Part II at agent-individual scale,
now extended to a 100+ named-individual cohort
(Section~\ref{sec:stage2-v55}).

\item Stage 4's Memetic Dynamics (Section~\ref{subsec:memetic})
operationalises the population-level analogue of the strategy
selection mechanism that NRMO\_CORE applies at the individual
decision level.

\item Stage 4's Black Swan catalog (Section~\ref{subsec:black-swan})
materialises the long-tail risk distribution against which the True
Ruin constraint of Part II was originally designed.
\end{itemize}
